\chapter{Fibre}

\section*{Definition and Overview}
Fibre, also spelled fiber, is the indigestible part of plant foods that is essential for a healthy digestive system and for the prevention of many chronic diseases. Unlike other carbohydrates, fibre is not broken down into sugar, so it does not raise blood glucose; instead, it passes through the gut, adding bulk to stools, feeding beneficial bacteria, and slowing the absorption of nutrients. A diet rich in fibre protects against constipation, haemorrhoids, diverticular disease, bowel cancer, heart disease, and type 2 diabetes. Despite these well-established benefits, most people eat well below the recommended amount. This chapter explains the two types of fibre, the foods that provide it, and practical ways to eat more.

\section*{Epidemiology}
Australians eat far less fibre than recommended. Guidelines suggest adults consume 25--30 grams per day, but average intake is around 20--25 grams, and many people, particularly those eating highly processed diets, consume considerably less. Low fibre intake is linked to the high prevalence of constipation (affecting about 15--20\% of people), haemorrhoids, and diverticular disease, and it contributes to Australia's high rates of bowel cancer, one of the most common cancers in the country. Population studies consistently show that higher fibre intake is associated with lower rates of heart disease, stroke, diabetes, and bowel cancer, making fibre one of the most protective components of the diet.

\section*{Aetiology and Pathophysiology}
\begin{itemize}
    \item \textbf{Soluble fibre:} dissolves in water to form a gel; found in oats, barley, psyllium, legumes, and fruit; lowers cholesterol and steadies blood sugar
    \item \textbf{Insoluble fibre:} does not dissolve; found in wholegrains, wheat bran, nuts, and vegetable skins; adds bulk and prevents constipation
    \item \textbf{Resistant starch:} a form of starch that resists digestion and feeds gut bacteria, found in cooked and cooled potatoes, legumes, and some grains
    \item \textbf{Mechanism:} fibre increases stool bulk, softens stool, and speeds transit, preventing constipation and straining
    \item \textbf{Cholesterol:} soluble fibre binds bile acids and cholesterol in the gut, increasing their excretion and lowering blood cholesterol
    \item \textbf{Blood sugar:} the gel slows sugar absorption, producing smaller glucose rises after meals
    \item \textbf{Microbiome:} fibre is fermented by gut bacteria into short-chain fatty acids that nourish the bowel wall, reduce inflammation, and support immunity
    \item \textbf{Satiety and weight:} fibre increases fullness and reduces energy intake
\end{itemize}

\section*{Clinical Features}
\begin{itemize}
    \item Regular bowel habits: soft, easily passed stools, typically between three times a day and three times a week
    \item Reduced constipation, straining, and haemorrhoid symptoms with adequate intake
    \item Better post-meal fullness and appetite control
    \item Steadier blood sugar levels in people with diabetes
    \item Lower LDL cholesterol on blood tests
    \item Temporary wind, bloating, and cramps when fibre is increased too quickly
    \item In people with irritable bowel syndrome, certain fibres may trigger symptoms, so intake is individualised
\end{itemize}

\section*{Differential Diagnosis}
Bowel symptoms require consideration of causes beyond fibre.
\begin{itemize}
    \item Constipation from inadequate fluid, inactivity, medicines, or medical conditions
    \item Irritable bowel syndrome, in which fibre tolerance varies
    \item Coeliac disease and other malabsorption conditions
    \item Bowel cancer, which must be considered with changed bowel habit, blood, or weight loss
    \item Diverticular disease, linked to long-term low fibre intake
    \item Normal variation in bowel frequency, which is wide
\end{itemize}

\section*{When to Seek Medical Care}
Most people can increase fibre without medical advice. A doctor should be consulted for persistent constipation, rectal bleeding, weight loss, or a change in bowel habit lasting more than a few weeks, and for advice tailored to irritable bowel syndrome or diabetes.

\textbf{Call 000 (or local emergency services) for:}
\begin{itemize}
    \item Severe abdominal pain with vomiting and the inability to pass wind or stool
    \item Sudden severe abdominal pain
    \item Significant bleeding from the rectum or collapse
    \item Any emergency symptoms
\end{itemize}

\section*{Diagnosis and Investigations}
\begin{itemize}
    \item \textbf{Dietary assessment:} review of usual food intake against the fibre target
    \item \textbf{Symptom history:} bowel frequency, stool consistency, straining, pain, and blood
    \item \textbf{Investigations where indicated:} blood tests, and colonoscopy for concerning symptoms or family history
    \item \textbf{No routine testing:} for most healthy people, increasing fibre is a dietary matter
\end{itemize}

\section*{Management}
\begin{itemize}
    \item \textbf{Wholegrains:} oats, wholegrain bread, brown rice, and wholegrain breakfast cereals
    \item \textbf{Legumes:} lentils, chickpeas, beans, and peas in meals several times a week
    \item \textbf{Vegetables and fruit:} five or more serves daily, with edible skins left on
    \item \textbf{Nuts and seeds:} a handful daily, including chia and flaxseed
    \item \textbf{Gradual increase:} add a serve of high-fibre food every few days, with fluids, to allow the gut to adapt
    \item \textbf{Supplements:} psyllium or other fibre supplements where dietary intake is insufficient
    \item \textbf{Fluids:} drink adequate water, since fibre without fluid can worsen constipation
    \item \textbf{Individualised plans:} for irritable bowel syndrome, coeliac disease, and diabetes, seek dietitian advice
\end{itemize}

\section*{Complications}
\begin{itemize}
    \item Constipation, haemorrhoids, and anal fissures from chronic low intake
    \item Diverticular disease
    \item Elevated risk of bowel cancer, heart disease, and diabetes
    \item Bloating and discomfort with rapid increases
    \item Rarely, bowel obstruction in people with strictures who consume very large amounts of insoluble fibre
\end{itemize}

\section*{Prognosis and Outcomes}
Fibre is one of the best-evidenced components of a healthy diet. Increasing intake improves bowel function within days to weeks, lowers cholesterol and blood sugar over months, and is associated with substantial reductions in chronic disease risk over a lifetime. The World Health Organization estimates that simply meeting fibre targets would prevent a significant share of bowel cancers and cardiovascular deaths. Because the change is dietary and gradual, it is achievable for most people and the benefits are durable.

\section*{Prevention and Self-Care}
\begin{itemize}
    \item Work towards 25--30 grams of fibre a day
    \item Make half your grain foods wholegrain
    \item Add legumes to soups, salads, and stews
    \item Eat vegetables at most meals and fruit as snacks
    \item Increase fibre gradually over weeks to avoid wind
    \item Drink plenty of water
    \item Use food labels to choose higher-fibre options
\end{itemize}

\section*{Sources and Further Reading}
The following reputable sources provide further information for healthcare students and patients seeking reliable, current advice about fibre.
\begin{itemize}
    \item National Health and Medical Research Council. \textit{Australian Dietary Guidelines.} https://www.eatforhealth.gov.au/
    \item Healthdirect Australia. \textit{Dietary fibre.} https://www.healthdirect.gov.au/fibre
    \item Dietitians Australia. \textit{Fibre.} https://dietitiansaustralia.org.au/
    \item Better Health Channel. \textit{Fibre.} https://www.betterhealth.vic.gov.au/
    \item NHS. \textit{How to get more fibre.} https://www.nhs.uk/live-well/eat-well/
    \item World Health Organization. \textit{Healthy diet.} https://www.who.int/
\end{itemize}
